\documentclass[12pt,a4paper]{ctexart}

% === Packages ===
\usepackage{geometry}
\geometry{left=2.5cm,right=2.5cm,top=2.5cm,bottom=2.5cm,headheight=15pt}
\usepackage{graphicx}
\usepackage{hyperref}
\usepackage{listings}
\usepackage{xcolor}
\usepackage{booktabs}
\usepackage{tabularx}
\usepackage{enumitem}
\usepackage{titlesec}
\usepackage{fancyhdr}
\usepackage{caption}
\usepackage{float}
\renewcommand{\tabularxcolumn}[1]{m{#1}}

% === Code style ===
\lstset{
  basicstyle=\ttfamily\small,
  breaklines=true,
  frame=single,
  numbers=left,
  numberstyle=\tiny,
  keywordstyle=\color{blue},
  commentstyle=\color{gray},
  stringstyle=\color{red},
  showstringspaces=false,
  tabsize=2,
}
% Rust language definition for listings (MiKTeX doesn't ship rust lang by default)
\lstdefinelanguage{Rust}{
  keywords={as,async,await,break,const,continue,crate,dyn,else,enum,extern,false,fn,for,if,impl,in,let,loop,match,mod,move,mut,pub,ref,return,self,Self,static,struct,super,trait,true,type,unsafe,use,where,while},
  sensitive=true,
  morecomment=[l]{//},
  morecomment=[s]{/*}{*/},
  morestring=[b]",
  morestring=[b]',
}
\lstdefinelanguage{JavaScript}{
  keywords={const,let,var,function,return,if,else,for,while,true,false,new,typeof,arguments,this,window,document,console},
  sensitive=true,
  morecomment=[l]{//},
  morecomment=[s]{/*}{*/},
  morestring=[b]",
  morestring=[b]',
  morestring=[b]`,
}

% === Header / Footer ===
\pagestyle{fancy}
\fancyhf{}
\fancyhead[L]{Pake Plus --- 实验报告}
\fancyhead[R]{\thepage}

% === Title Page ===
\title{
  \vspace{3cm}
  {\Huge \textbf{Pake Plus}}\\[0.5cm]
  {\Large 基于 Pake 的网页桌面应用增强工具}\\[1cm]
  {\large 实验报告}
  \vspace{2cm}
}
\author{
  \begin{tabular}{ll}
  成员 A & 学号 XXXXXX（广告/跟踪拦截）\\
  成员 B & 学号 XXXXXX（离线缓存代理）\\
  成员 C & 学号 XXXXXX（剪贴板管理）\\
  成员 D & 学号 XXXXXX（设置面板 + 数据迁移 + 诊断 + 视频）\\
  \end{tabular}
}
\date{\today}

\begin{document}

\maketitle
\thispagestyle{empty}
\newpage

% === Abstract ===
\begin{center}
  \Large\textbf{摘要}
\end{center}
\vspace{0.5cm}

Pake Plus 是基于开源项目 Pake（Rust + Tauri v2）的网页桌面应用增强工具。
原 Pake 可将任意网页打包为约 5MB 的轻量级桌面应用（仅为 Electron 的 1/20），
但打包产物仅具备基础窗口功能。

本项目在 Pake 基础上新增四大模块：
\textbf{广告/跟踪拦截}（EasyList 规则引擎，双重过滤）、
\textbf{离线缓存代理}（HTTP 层透明缓存，LRU 淘汰）、
\textbf{剪贴板管理}（系统级 FFI 监听，全文搜索）和
\textbf{设置面板}（可视化配置、数据迁移、系统诊断、版本回滚）。

技术栈为 Rust + Tauri v2 IPC + 原生 JavaScript，四个模块通过共享 JSON 配置文件实现解耦协作。
Rust 后端负责系统级操作（文件 I/O、ZIP 打包、系统信息采集、OS FFI），
JavaScript 前端通过纯 DOM 操作构建用户界面。
项目总计 67 个 Rust 单元测试、297 个 JavaScript 测试，全部通过，0 warnings。

\vspace{0.5cm}
\noindent\textbf{关键词}：Rust；Tauri；桌面应用；广告拦截；离线缓存；剪贴板管理；WebView

\newpage

% === TOC ===
\tableofcontents
\newpage

% ================================================================
% Chapter 1
% ================================================================
\section{项目名称与团队介绍}

\subsection{项目简介}

Pake Plus 是基于开源项目 Pake（Rust + Tauri v2）的网页桌面应用增强工具。原 Pake 项目可将任意网页打包为轻量级原生桌面应用（约 5MB），但功能仅限基础窗口管理和系统托盘。Pake Plus 在此基础上新增四个独立功能模块，使其更接近一个成熟的桌面浏览器应用。

\subsection{团队成员与分工}

\begin{table}[h]
\centering
\small
\begin{tabularx}{\textwidth}{c>{\centering\arraybackslash}p{3.5cm}X}
\toprule
成员   & 负责模块        & 核心工作 \\
\midrule
成员 A & 广告/跟踪拦截   & EasyList 规则引擎、网络请求拦截、DOM 元素隐藏、自定义规则、拦截统计 \\
成员 B & 离线缓存代理   & HTTP 代理拦截、磁盘缓存管理、LRU 淘汰算法、离线模式 \\
成员 C & 剪贴板管理     & 系统剪贴板监听（FFI）、历史记录面板、全文搜索、自动清理 \\
成员 D & 设置面板 + 数据 & 可视化配置页面、数据导出/导入（ZIP）、系统诊断、视频制作 \\
\bottomrule
\end{tabularx}
\end{table}

% ================================================================
% Chapter 2
% ================================================================
\section{项目背景与目标}

\subsection{选题原因}

Pake 是一个优秀的轻量级桌面应用打包工具，利用系统 WebView 代替 Electron 的 Chromium 内核，
打包体积仅约 5MB（约 Electron 的 1/20）。然而，原项目在功能层面存在明显不足：
打包后的应用只是一个"空壳浏览器"，缺少现代桌面浏览器应具备的以下能力：

\begin{enumerate}[label=(\arabic*)]
  \item 无广告过滤能力 --- 网页中的广告和跟踪脚本原样加载
  \item 无离线缓存 --- 断网后页面完全无法访问
  \item 无剪贴板历史 --- 依赖系统剪贴板，复制内容无法回溯
  \item 无可视化配置 --- 所有参数在打包时通过 CLI 一次性写入，运行时不可修改
\end{enumerate}

\subsection{改进目标}

本项目旨在解决上述四个痛点，在不显著增加打包体积的前提下，通过 Rust 后端 + Tauri IPC 框架，
为 Pake 打包的应用增加广告拦截、离线缓存、剪贴板管理和可视化设置四大功能模块。

\subsection{与原项目 Pake 的区别}

\begin{table}[h]
\centering
\small
\begin{tabularx}{\textwidth}{lXX}
\toprule
对比维度 & 原 Pake & Pake Plus \\
\midrule
广告处理   & 无任何广告过滤能力 & EasyList 规则引擎，网络层拦截 + DOM 元素隐藏双重过滤 \\
离线能力   & 完全依赖网络 & HTTP 层代理缓存，已访问页面断网仍可浏览 \\
剪贴板能力 & 无历史记录 & 系统剪贴板监听，全文搜索 + 一键复用 \\
功能配置   & 打包时 CLI 一次性配置 & 运行时可视化面板，修改即时生效并持久化 \\
数据迁移   & 不支持 & 一键导出 ZIP，换设备导入恢复 \\
系统诊断   & 无 & 编译信息 + 系统信息，一键复制报告 \\
\bottomrule
\end{tabularx}
\end{table}

% ================================================================
% Chapter 3
% ================================================================
\section{系统设计与整体架构}

\subsection{技术栈}

\begin{itemize}
  \item \textbf{后端}：Rust（Tauri v2 框架）
  \item \textbf{前端}：TypeScript（CLI）+ JavaScript（WebView 注入）
  \item \textbf{桌面壳}：系统 WebView（Windows: WebView2，macOS: WKWebView，Linux: WebKitGTK）
  \item \textbf{IPC}：Tauri Command（Rust ↔ JavaScript 双向调用）
  \item \textbf{关键 Crates}：zip 2.2、sysinfo 0.33、arboard 3.4、chrono 0.4、built 0.7、clipboard-rs 0.3、rusqlite 0.32、sha2 0.10、regex 1、url 2.5、serde / serde\_json
\end{itemize}

\subsection{整体架构图}

\begin{figure}[h]
\centering
% TODO: 插入架构图
% \includegraphics[width=\textwidth]{architecture.png}
\caption{Pake Plus 系统架构（待插入图片）}
\end{figure}

\subsection{模块关系}

四个模块在功能层面完全独立，互不依赖：

\begin{itemize}
  \item A 的广告拦截不需要 B/C/D 的任何功能即可运行
  \item B 的离线缓存不需要 A/C/D 的任何功能即可运行
  \item C 的剪贴板管理不需要 A/B/D 的任何功能即可运行
  \item D 的设置面板通过读/写共享配置文件管理 A/B/C 模块，但不影响各模块独立运行
\end{itemize}

四个模块通过 D 定义的 \texttt{AppSettings} 数据结构进行配置共享，
通过 Tauri 的 \texttt{\#[command]} IPC 机制暴露接口给前端调用。

\subsection{文件结构}

\begin{lstlisting}[caption=Pake Plus 代码结构]
src-tauri/src/
├── adblock/              # [A] 广告拦截模块
│   ├── mod.rs            # 模块入口 + AdblockState
│   ├── engine.rs         # URL 匹配引擎（域名/正则）
│   └── rules.rs          # EasyList 规则解析器
├── app/
│   ├── clipboard/          # [C] 剪贴板管理模块
│   │   ├── mod.rs          # 模块入口
│   │   ├── monitor.rs      # 系统剪贴板监听（FFI）
│   │   ├── store.rs        # SQLite 存储与检索
│   │   ├── filter.rs       # 隐私过滤（密码/银行卡/Luhn）
│   │   ├── panel.rs        # 剪贴板历史面板窗口
│   │   ├── commands.rs     # IPC 命令
│   │   ├── settings.rs     # 剪贴板专用设置
│   │   ├── cleanup.rs      # 自动清理
│   │   └── source.rs       # 来源应用识别
│   ├── settings/          # [D] 设置面板模块
│   │   ├── mod.rs         # 模块入口 + 重导出
│   │   ├── types.rs       # 数据结构 + Theme/Language 枚举
│   │   ├── traits.rs      # ModuleSettings trait
│   │   ├── io.rs          # 文件读写 + 版本备份
│   │   ├── health.rs      # 启动健康检查
│   │   ├── diagnostics.rs # 诊断信息采集
│   │   ├── commands.rs    # IPC 命令
│   │   └── tests.rs       # 单元测试
│   ├── setup.rs           # 托盘菜单 + 全局快捷键
│   ├── window.rs          # 窗口创建 + JS 注入
│   └── ...
├── inject/
│   ├── custom.js          # [D] 设置面板侧边栏（齿轮按钮触发）
│   ├── settings.js        # [C] 剪贴板测试面板
│   ├── adblock.js         # [A] fetch/XHR 拦截 + DOM 隐藏
│   ├── event.js           # 剪贴板快捷键 + 事件处理
│   └── ...
└── lib.rs                 # 应用入口 + 命令注册
\end{lstlisting}

\subsection{技术亮点}

\begin{table}[htbp]
\centering
\small
\begin{tabularx}{\textwidth}{lXc}
\toprule
亮点 & 说明 & 涉及模块 \\
\midrule
\textbf{Rust 枚举类型安全} & Theme/Language 从 String 重构为 enum，非法值编译期拒绝，\texttt{\#[serde(alias)]} 向后兼容旧格式 & D \\
\textbf{ModuleSettings trait 多态} & 每个模块独立实现 \texttt{validate() / summary() / stats()}，IPC 统一调用 & D \\
\textbf{跨平台剪贴板监听} & Windows、macOS 与 X11 使用 \texttt{clipboard\_rs} 原生事件监听，纯 Wayland 使用 \texttt{wl\_clipboard\_rs} 低频轮询兜底，独立线程不阻塞主界面 & C \\
\textbf{双重哈希去重与回写抑制} & 基于 SHA-256 的 \texttt{last\_hash} 与 \texttt{expected\_hash} 双重判断，过滤系统重复通知，并避免从历史面板复制时再次入库 & C \\
\textbf{隐私过滤与来源识别} & 识别 Windows、macOS、X11 来源应用，支持忽略应用、疑似密码检测和银行卡号 Luhn 校验，敏感内容在入库前即被拦截 & C \\
\textbf{SQLite 历史检索与自动清理} & SQLite WAL 持久化、唯一哈希索引与 UPSERT 去重，支持多关键词 AND 模糊搜索，并按保留天数和容量上限定时清理 & C \\
\textbf{HTTP 层透明代理缓存} & 基于 Tauri 资源拦截，异步磁盘 I/O，LRU 淘汰算法 & B \\
\textbf{EasyList 规则引擎} & 网络请求拦截 + DOM 元素隐藏双重过滤，支持自定义规则 & A \\
\textbf{JavaScript 请求代理} & 代理 window.fetch / XMLHttpRequest，匹配规则返回空响应，MutationObserver 动态隐藏广告元素 & A \\
\textbf{运行时规则同步} & 设置面板保存自定义规则后通过 IPC 同步更新广告引擎，即时生效 & A \\
\textbf{设置版本备份与自动恢复} & 保存前轮转 5 个历史版本，启动时健康检查自动从备份修复损坏文件 & D \\
\textbf{模块化文件结构} & 800 行单文件拆为 8 个模块（types / traits / io / health / diagnostics / commands / tests） & D \\
\textbf{数据可移植性} & 动态扫描数据目录打包 ZIP，导入前预览确认，先备份后覆盖 & D \\
\bottomrule
\end{tabularx}
\caption{技术亮点汇总}
\end{table}

% ================================================================
% Chapter 4 (by member)
% ================================================================
\section{功能模块实现}

% ---------- 4.1 A ----------
\subsection{广告拦截模块（成员 A）}

\subsubsection{设计思路}

广告拦截模块采用``规则引擎 + 前端拦截''的双层架构。Rust 后端负责解析 EasyList 格式的规则文本，将其分类为域名过滤规则、正则表达式规则和 CSS 隐藏选择器三类，通过 Tauri 的 \texttt{initialization\_script} 机制序列化为 JSON 注入到页面中。前端 JavaScript（\texttt{adblock.js}）在页面加载时读取 \texttt{window.pakeAdblock} 配置，随后代理浏览器原生的 \texttt{fetch} 和 \texttt{XMLHttpRequest} 对象——对匹配黑名单的网络请求直接返回 HTTP 204 空响应，阻止广告资源加载；同时将 CSS 隐藏选择器注入页面 \texttt{<style>} 标签，并通过 MutationObserver 监听 DOM 变化，确保后续动态插入的广告元素也被实时隐藏。

模块的构建时配置通过 \texttt{--block-ads} CLI 选项启用，通过 \texttt{--adblock-rules <path>} 加载用户自定义规则文件（每行一条，支持 \texttt{||domain\textasciicircum} 域名匹配和 \texttt{\#\#.selector} CSS 隐藏两种格式）。运行时，设置面板的 Adblock 标签提供 ON/OFF 开关和折叠式规则编辑器——开关通过全局变量 \texttt{window.\_\_pakeAdblockRunning} 实时控制 \texttt{adblock.js} 的拦截行为（无需页面刷新）；自定义规则保存时，Rust 后端的 \texttt{save\_settings} 命令自动对比新旧规则并调用引擎的 \texttt{add\_custom\_rule / remove\_custom\_rule} 完成同步。拦截计数器通过托盘 tooltip 显示，设置面板的拦截统计区域通过 \texttt{adblock\_get\_stats} IPC 命令获取实时数据。

技术难点在于：初始化脚本的执行时机早于页面内容加载，需要确保 \texttt{AdblockState} 在窗口构建前完成初始化（否则 \texttt{try\_state} 返回 None 导致 \texttt{adblock.js} 不被注入）；Tauri v2 的 ACL 权限系统默认拒绝所有自定义 IPC 命令，需要将所有命令统一配置在权限文件中；CSS 隐藏选择器需要兼容 EasyList 的多种语法变体（\texttt{\#\#}、\texttt{\#\#\#}、\texttt{\#@\#} 例外规则等）。

\subsubsection{核心流程}

\begin{sloppypar}
\begin{enumerate}[label=(\arabic*)]
  \item \textbf{构建时规则加载。}
    用户执行 CLI 命令 \texttt{node dist/cli.js <url> --block-ads --adblock-rules ./my-rules.txt} 时，CLI 将 \texttt{block\_ads} 和 \texttt{adblock\_rules} 写入 \texttt{.pake/pake.json}。Rust 编译期间通过 \texttt{include\_str!} 读取内置规则文件 \texttt{easylist-mini.txt}（包含常用广告和跟踪域名），与用户自定义规则合并后初始化 \texttt{AdblockEngine}。

  \item \textbf{规则引擎初始化与状态管理。}
    应用启动时，\texttt{run\_app()} 首先将 \texttt{pake\_config.block\_ads} 和 \texttt{adblock\_rules} 传递给 \texttt{AdblockState::new}。该函数调用 \texttt{AdblockEngine::from\_rules\_text} 解析所有规则文本：\texttt{parse\_rules\_text} 按行遍历，通过正则匹配识别规则类型（\texttt{||domain\textasciicircum} → 域名规则、\texttt{/regex/} → 正则规则、\texttt{\#\#selector} → CSS 隐藏选择器），结果存入 \texttt{EngineRules} 结构体。初始化完成后通过 \texttt{app.manage()} 将状态注册到 Tauri 状态管理器。

  \item \textbf{规则导出与前端注入。}
    窗口构建函数 \texttt{build\_window} 检查 \texttt{config.block\_ads} 标志，若启用则通过 \texttt{app.try\_state::<AdblockState>()} 获取引擎实例，调用 \texttt{export\_for\_injection(None)} 将内存中的域名列表（\texttt{Vec<String>}）、正则表达式（\texttt{Vec<String>}）、CSS 选择器（\texttt{Vec<String>}）序列化为 JSON，以 \texttt{window.pakeAdblock = \{enabled:true, domains:[...], regexes:[...], cosmetic\_selectors:[...]\}} 格式注入为初始化脚本。随后注入 \texttt{adblock.js} 文件。

  \item \textbf{前端网络请求拦截。}
    \texttt{adblock.js} 执行时首先检查 \texttt{window.pakeAdblock} 是否存在且 \texttt{enabled} 为 true。若启用，保存原始 \texttt{fetch} 和 \texttt{XMLHttpRequest.open} 引用，然后替换为包裹函数。包裹函数调用 \texttt{shouldBlockUrl()} 对请求 URL 进行匹配——域名规则直接比较 hostname（支持精确匹配和子域名匹配），正则规则对整个 URL 进行 \texttt{RegExp.test()}。匹配成功则调用 \texttt{reportBlock()} 递增计数器，然后返回 HTTP 204 空响应或替换 \texttt{XMLHttpRequest} 为 \texttt{about:blank}，阻止实际网络请求。

  \item \textbf{前端 CSS 隐藏与 DOM 监听。}
    \texttt{applyCosmeticRules()} 遍历 \texttt{cosmetic\_selectors} 列表，删除 \texttt{\#\#} 和 \texttt{\#\#\#} 前缀，将剩余选择器拼接为 \texttt{\{display:none!important;visibility:hidden!important\}} 样式，注入到页面 \texttt{<style id="pake-adblock-style">} 标签中。同时创建 \texttt{MutationObserver} 监听 \texttt{documentElement} 的 \texttt{childList} 和 \texttt{subtree} 变化，每当有新元素插入时重新应用 CSS 规则，确保动态加载的广告也被隐藏。

  \item \textbf{运行时开关控制。}
    设置面板的 Adblock 标签页渲染 ON/OFF 开关（\texttt{id=\_\_ps\_adblk\_\_}），点击时通过 \texttt{tg()} 函数的 onclick 处理器设置全局变量 \texttt{window.\_\_pakeAdblockRunning}。\texttt{adblock.js} 的 \texttt{shouldBlockUrl()} 和 \texttt{reportBlock()} 在每次执行前通过 \texttt{isAdblockEnabled()} 检查该标志——若为 \texttt{false}，所有拦截逻辑立即返回，请求正常通行。

  \item \textbf{拦截统计展示。}
    设置面板的拦截统计区域通过 \texttt{loadAdblockStats()} 调用 \texttt{adblock\_get\_stats} IPC 命令获取实时数据（拦截次数、规则数量）。托盘菜单的 tooltip 通过 \texttt{adblock::update\_tray\_block\_count} 更新为``已拦截 N 个请求''。前端拦截计数器每拦截一次递增，同时调用 \texttt{adblock\_report\_blocked} IPC 通知后端同步更新。

  \item \textbf{自定义规则生命周期。}
    用户点击``✏ 编辑自定义规则''展开文本区，输入规则后点击 Save。\texttt{doSave()} 收集 \texttt{\_\_ps\_rules\_\_} 文本区的值构造 \texttt{AppSettings.adblock.custom\_rules}，调用 \texttt{save\_settings} IPC。Rust 后端的 \texttt{save\_settings} 命令在写入 JSON 文件后，检测广告规则是否变化——若有差异，依次调用 \texttt{remove\_custom\_rule} 清除旧规则、\texttt{add\_custom\_rule} 添加新规则，每个规则再次经 \texttt{parse\_custom\_rule} 校验后写入引擎的 \texttt{custom\_lines} 列表，并重新合并生成网络规则和 CSS 规则。
\end{enumerate}
\end{sloppypar}

\subsubsection{关键代码}

\textbf{1. 规则解析与分类}

\begin{lstlisting}[language=Rust, caption=EasyList 规则解析器]
pub struct ParsedRules {
    pub network: Vec<NetworkRule>,
    pub cosmetic: Vec<CosmeticRule>,
}
pub fn parse_rules_text(text: &str) -> ParsedRules {
    let mut network = Vec::new(); let mut cosmetic = Vec::new();
    for line in text.lines() {
        let trimmed = line.trim();
        if trimmed.is_empty() || trimmed.starts_with('!') { continue; }
        if trimmed.starts_with("##") { cosmetic.push(CosmeticRule { ... }); }
        else if trimmed.starts_with("||") { network.push(NetworkRule::Domain(...)); }
        else if trimmed.starts_with('/') { network.push(NetworkRule::UrlRegex(...)); }
    }
    ParsedRules { network, cosmetic }
}
\end{lstlisting}

引擎启动时通过 \texttt{parse\_rules\_text} 将 EasyList 规则文本解析为结构化数据。域名规则（\texttt{||domain\textasciicircum}）存储为纯字符串用于 hostname 匹配；正则规则（\texttt{/pattern/}）编译为 \texttt{Regex} 对象并用 \texttt{Arc} 包装以支持跨线程共享；CSS 隐藏规则（\texttt{\#\#selector}）存储为选择器字符串。

\textbf{2. 配置注入与初始化顺序}

\begin{lstlisting}[language=Rust, caption=AdblockState 初始化和窗口注入]
// lib.rs setup closure — order matters!
let adblock_state = adblock::AdblockState::new(true, &adblock_rules);
app.manage(adblock_state);  // MUST be before set_window

let window = set_window(app.app_handle(), &pake_config, &tauri_config)?;
// Inside build_window():
if config.block_ads {
    if let Some(state) = app.try_state::<adblock::AdblockState>() {
        let export = state.engine.export_for_injection(None);
        window_builder = window_builder
            .initialization_script(&format!("window.pakeAdblock={};", json!(export)))
            .initialization_script(include_str!("../inject/adblock.js"));
    }
}
\end{lstlisting}

\texttt{app.manage()} 必须在 \texttt{set\_window} 之前调用，否则 \texttt{build\_window} 中的 \texttt{try\_state} 返回 None，导致 \texttt{adblock.js} 不被注入。这是最常见的初始化顺序问题。

\textbf{3. 前端拦截代理与运行时开关}

\begin{lstlisting}[language=JavaScript, caption=adblock.js 核心逻辑]
const originalFetch = window.fetch;
window.fetch = function (input, init) {
    if (!isAdblockEnabled()) return originalFetch.apply(this, arguments);
    const url = typeof input === 'string' ? input : input?.url;
    if (shouldBlockUrl(url)) { reportBlock(); return new Response('', {status:204}); }
    return originalFetch.apply(this, arguments);
};

function applyCosmeticRules() {
    const css = config.cosmetic_selectors.map(s =>
        s.replace(/^##/, '') + '{display:none!important}').join('\n');
    let style = document.getElementById('pake-adblock-style');
    if (!style) { style = document.createElement('style'); style.id = 'pake-adblock-style';
        (document.head || document.documentElement).appendChild(style); }
    style.textContent = css;
}
const observer = new MutationObserver(() => applyCosmeticRules());
observer.observe(document.documentElement, { childList: true, subtree: true });
\end{lstlisting}

\textbf{4. 设置面板到引擎的规则同步}

\begin{lstlisting}[language=Rust, caption=save\_settings 中的规则同步]
// In settings/commands.rs save_settings()
if let Some(state) = app.try_state::<crate::adblock::AdblockState>() {
    let current = state.engine.custom_rules();
    if current != settings.adblock.custom_rules {
        for rule in &current { state.engine.remove_custom_rule(rule); }
        for rule in &settings.adblock.custom_rules {
            let _ = state.engine.add_custom_rule(rule);
        }
        eprintln!("[Pake] adblock rules synced: {} rules", settings.adblock.custom_rules.len());
    }
}
\end{lstlisting}

\subsubsection{效果截图}

% TODO: 成员 A 插入实际运行截图
\begin{itemize}
  \item CLI 构建命令与 \texttt{--block-ads} 选项输出
  \item 左下角 Pake Plus 按钮与页面正常运行状态
  \item 设置面板 Adblock 标签——ON/OFF 开关、ON 标签、拦截统计数字、✏ 编辑自定义规则
  \item 自定义规则编辑区展开状态（文本区 + 格式提示）
  \item 托盘菜单 tooltip 显示``已拦截 N 个请求''
\end{itemize}

\subsubsection{遇到的问题与解决方案}

\begin{enumerate}[label=(\arabic*)]
  \item \textbf{AdblockState 初始化时机晚于窗口创建，导致 adblock.js 不注入}
    \\→ 解决：将 \texttt{app.manage(adblock\_state)} 移到 \texttt{set\_window} 调用之前。初始化顺序为：创建 AdblockState → manage → set\_window → build\_window 中 try\_state 获取状态 → 注入脚本。若顺序错误，\texttt{try\_state} 返回 None 且无任何错误提示，通过 \texttt{eprintln!} 日志定位。

  \item \textbf{设置面板自定义规则保存后不生效——JSON 存储与引擎内存脱节}
    \\→ 解决：在 \texttt{save\_settings} IPC 命令末尾增加规则同步逻辑。对比当前引擎中的 \texttt{custom\_rules()} 与新设置中的规则列表，先清除所有旧规则再逐条添加新规则。每个规则添加前通过 \texttt{parse\_custom\_rule} 校验格式，非法规则不进入引擎。

  \item \textbf{Tauri v2 ACL 拦截所有自定义 IPC 命令（adblock\_get\_stats 等无法调用）}
    \\→ 解决：将 adblock 的命令名写入 \texttt{permissions/settings.toml} 的 \texttt{commands.allow} 列表，确保与 \texttt{capabilities/default.json} 中的权限标识符对应。Tauri v2 要求每个自定义命令显式授权，分散在多个权限文件容易遗漏。

  \item \textbf{开关切换后拦截不停止——adblock.js 在加载时已固化状态}
    \\→ 解决：在 \texttt{adblock.js} 中增加 \texttt{isAdblockEnabled()} 函数，在每次 \texttt{shouldBlockUrl} 和 \texttt{reportBlock} 执行前动态检查 \texttt{window.\_\_pakeAdblockRunning} 全局变量。设置面板的开关 onclick 和 \texttt{fillAdblock} 加载时同步更新该变量。

  \item \textbf{多个齿轮按钮并存（旧右下角 FAB + 新左下角按钮）}
    \\→ 解决：定位 \texttt{custom.js} 中两处按钮创建代码——旧 FAB（\texttt{\_\_ps\_fab}）和旧 \texttt{\_\_pgb}——逐一移除，统一使用左下角白色卡片样式按钮。

  \item \textbf{自定义规则文本区默认不可见——用户不知道如何编辑}
    \\→ 解决：改为折叠面板设计——默认显示``✏ 编辑自定义规则''按钮（虚线边框），点击后展开文本区并切换为``✏ 收起''（蓝色实线边框），提供清晰的展开/收起视觉反馈。
\end{enumerate}

% ---------- 4.2 B ----------
\subsection{离线缓存代理模块（成员 B）}

\subsubsection{设计思路}

\begin{itemize}
  \item \textbf{目标}：在 HTTP 层透明代理请求，将响应缓存到磁盘，支持离线访问
  \item \textbf{Rust 技术}：异步 I/O（tokio）、资源拦截、LRU 缓存淘汰、HTTP 头解析
  \item \textbf{核心挑战}：大量并发请求的缓存写入吞吐量；缓存过期策略的准确性；断网判断
\end{itemize}

% TODO: 成员 B 完善设计思路（300-500 字）

\subsubsection{核心流程}

\begin{enumerate}[label=(\arabic*)]
  \item 注册 Tauri HTTP 拦截器，拦截所有 GET 请求
  \item 请求到达时查本地缓存索引，命中且未过期直接返回缓存数据
  \item 未命中或已过期则放行请求，获取响应后异步写入磁盘
  \item 维护缓存文件索引（URL 哈希 → 文件路径 → 大小 → 最后访问时间）
  \item 缓存总量超过上限（默认 200MB）时触发 LRU 淘汰
  \item 断网时展示离线提示页，列出可访问的已缓存页面
\end{enumerate}

% TODO: 成员 B 完善流程描述

\subsubsection{关键代码}

\begin{lstlisting}[language=Rust, caption=离线缓存示例代码]
% TODO: 成员 B 插入核心代码片段（HTTP 拦截器、LRU 淘汰、缓存索引维护等）
% 建议 1-2 段关键代码，每段 5-20 行
\end{lstlisting}

\subsubsection{效果截图}

% TODO: 成员 B 插入截图
% 建议：1) 正常缓存命中统计  2) 断网后页面正常显示  3) 缓存管理设置页  4) 离线提示页

\subsubsection{遇到的问题与解决方案}

% TODO: 成员 B 填写

% ---------- 4.3 C ----------
\subsection{剪贴板管理模块（成员 C）}

\subsubsection{设计思路}

剪贴板管理模块的目标是在不打断用户原有复制、粘贴习惯的前提下，为 Pake 应用补充可检索、
可复用的剪贴板历史。模块采用“系统监听层—隐私过滤层—持久化层—交互层”的分层结构：
监听层通过 \texttt{clipboard\_rs} 对接 Windows、macOS 和 X11 的原生剪贴板事件，
在纯 Wayland 环境下使用 300 ms 低频轮询作为兼容方案；监听任务运行在独立线程中，
避免阻塞 Tauri 主线程。每次检测到文本变化后先计算 SHA-256 摘要，并结合
\texttt{last\_hash} 与 \texttt{expected\_hash} 消除连续重复内容以及“从历史记录复制”造成的回写。

通过初步去重的文本还需经过隐私过滤：空文本、过短或过长文本、疑似密码、通过 Luhn 校验的
银行卡号，以及用户指定应用产生的内容均可跳过。有效记录存入 SQLite 数据库，数据库启用 WAL
模式并为内容哈希和创建时间建立索引；相同内容再次出现时通过 UPSERT 更新时间，从而兼顾去重
和最近使用排序。前端使用独立的 Tauri Webview 面板，通过 IPC 完成列表加载、关键词搜索、复制、
删除、清空和设置更新。该设计使监听、存储和界面相互解耦，同时兼顾跨平台能力、性能与隐私安全。

\subsubsection{核心流程}

\begin{sloppypar}
\begin{enumerate}[label=(\arabic*)]
  \item \textbf{初始化配置与运行状态。}
    应用启动时，\path{init_clipboard_state} 首先通过 \path{get_data_dir}
    获取应用数据目录，并定位 \path{clipboard-settings.json}。随后调用
    \path{ClipboardSettings::load} 读取持久化设置；若文件不存在或 JSON 损坏，
    则回退到构建参数 \path{clipboard} 与 \path{clipboard_max} 所确定的默认值。
    \path{normalized} 会再次执行构建开关校验，并将最大记录数限制在 500--5000 条之间，
    防止运行时配置绕过打包阶段的能力限制。设置保存采用“先写临时文件，再重命名替换”的方式，
    以降低程序异常退出造成配置文件只写入一半的风险。

  \item \textbf{创建 SQLite 历史数据库。}
    仅当剪贴板功能启用时，后端才调用 \path{ClipboardStore::open}，在
    \path{clipboard/history.db} 中创建数据库。数据库表保存编号、文本内容、SHA-256 摘要、
    创建时间和来源应用五类信息，并启用 WAL 日志模式。程序为 \path{content_hash}
    建立唯一索引，为 \path{created_at} 建立普通索引；启动时还会检测旧数据库是否缺少
    \path{source_app} 字段，并通过 \texttt{ALTER TABLE} 自动迁移。数据库打开后立即执行一次
    清理，避免过期数据在应用重新启动后继续保留。

  \item \textbf{选择与启动平台监听方式。}
    \path{start_monitor} 先判断当前 Linux 会话是否为纯 Wayland：当
    \path{WAYLAND_DISPLAY} 存在且 \path{DISPLAY} 为空时，启动
    \path{wl_clipboard_rs} 轮询线程，每 300 ms 读取一次文本；其他环境则通过
    \path{clipboard_rs} 创建 \path{ClipboardWatcherContext}，注册
    \path{HistoryHandler} 接收原生剪贴板变化事件。监听启动前会读取一次当前剪贴板并保存摘要，
    因而不会把应用启动前已经存在的内容误记为新记录。监听线程同时保存关闭通道或原子停止标志，
    在模块停用和应用退出时可以停止循环并 \path{join} 回收线程。

  \item \textbf{读取变化并执行双重去重。}
    系统通知到达后，\path{HistoryHandler::on_clipboard_change} 调用
    \path{get_text} 获取纯文本，并通过 \path{ClipboardStore::hash_content}
    计算 SHA-256。第一层去重检查共享的 \path{expected_hash}：若该摘要来自历史面板主动写回，
    则消费标记并直接返回；第二层检查监听器自己的 \path{last_hash}，过滤操作系统可能连续发送的
    相同通知。只有两次检查均未命中的文本才会进入 \path{process_text}，从源头减少过滤和数据库写入开销。

  \item \textbf{识别剪贴板来源应用。}
    去重完成后，后端调用 \path{current_application} 获取来源名称。Windows 优先读取
    \path{GetClipboardOwner}，若应用已释放所有权则回退到前台窗口，再根据进程编号查询可执行文件名；
    macOS 通过 \path{NSWorkspace} 获取最前端应用；X11 根据
    \path{_NET_ACTIVE_WINDOW} 和 \path{_NET_WM_PID} 查找进程名称。
    纯 Wayland 出于协议权限限制无法获知其他客户端的来源，因此允许该字段为空。

  \item \textbf{执行隐私与有效性过滤。}
    \path{process_text} 先从 \path{RwLock} 读取最新设置；模块已关闭时立即结束。
    随后构造 \path{FilterConfig} 并调用 \path{should_record}：空文本和少于 2 个字符的
    短文本默认跳过，超过 10000 个字符的内容拒绝记录；来源应用名称会去除路径及
    \path{.exe}/\path{.app} 后缀，再与忽略列表进行不区分大小写的比较；13--19 位数字候选值
    只有通过 Luhn 校验才按银行卡号过滤；6--30 位、无空白且同时包含字母和数字的内容按疑似密码处理。
    这些隐私开关可在本地历史面板中修改，过滤失败的内容不会进入数据库。

  \item \textbf{持久化记录并保持最近使用顺序。}
    合格文本交给 \path{insert_with_source}。该函数再次计算摘要并生成 UTC RFC 3339 时间，
    然后使用参数化 SQL 写入数据库。若唯一索引发现相同 \path{content_hash} 已存在，
    \texttt{ON CONFLICT DO UPDATE} 不新增行，而是刷新文本、时间和可用的来源应用；因此重复复制的
    内容只会移动到历史列表顶部，不会无限产生副本。查询列表时按 \texttt{created\_at DESC, id DESC}
    排序，默认返回最近 100 条，IPC 层把允许的返回数量限制在 1--500 条。

  \item \textbf{定期清理过期与超量数据。}
    每次插入结束后都会调用 \path{cleanup}，先删除早于保留期限的记录；默认保留 30 天。
    若剩余数量超过 \path{max_items}，则按创建时间升序选出最早记录，并以 500 条为批次淘汰，
    避免为每一条超限数据启动一次删除事务。除此之外，\path{start_cleanup} 还创建名为
    \path{pake-clipboard-cleanup} 的后台线程，通过 \path{recv_timeout} 每小时执行一次清理，
    即使长时间没有新的复制操作，过期记录也能被移除；停止信号到达后线程立即退出。

  \item \textbf{通过全局快捷键或设置面板打开历史面板。}
    功能启用后，\path{set_global_shortcut} 在 Windows/Linux 注册 Ctrl+Shift+V，在 macOS 注册
    Cmd+Shift+V。若它与应用唤醒快捷键相同，则剪贴板历史优先；回调还通过
    \path{Instant} 设置 300 ms 防抖。此外，设置面板的剪贴板标签页提供蓝色“打开剪贴板历史”按钮，
    点击后通过 Tauri 事件系统（\texttt{emit} → \texttt{listen}）通知 Rust 后端打开面板。
    触发后 \path{toggle_panel} 创建或复用
    \path{clipboard-panel} Webview：初始大小为 $420\times520$，最小为 $320\times360$，
    窗口带标题栏、居中显示、可缩放且显示在任务栏中。
    关闭按钮隐藏窗口而非销毁，面板已存在时只刷新内容并获取焦点。

  \item \textbf{加载列表并完成多关键词搜索。}
    面板显示后调用 \path{clipboard_list} 加载最近记录，同时读取
    \path{clipboard_stats} 展示“当前条数/容量上限”。用户输入搜索内容时，前端先等待
    120 ms，以减少连续输入造成的 IPC 请求；随后调用 \path{clipboard_search}。
    后端按空白拆分关键词，对 \texttt{\%}、\texttt{\_} 和反斜杠进行转义，再为每个词生成
    \texttt{LOWER(content) LIKE} 条件，并用 AND 连接。因此输入多个关键词时，记录必须同时包含
    所有词才会命中；所有值均通过 SQL 参数绑定传入，避免查询语义被特殊字符改变。

  \item \textbf{复用、删除和清空历史。}
    用户点击记录或“复制”按钮后，前端将记录编号传给 \path{clipboard_copy_item}；后端先按
    ID 读取完整文本，再由 \path{write_text_to_clipboard} 写回系统剪贴板。写入前保存
    \path{expected_hash}，写入失败时立即撤销，成功后若 1 秒内没有对应通知也会自动清除，
    防止标记长期残留。复制成功后面板保持打开，并在鼠标附近显示 3 秒提示。单条删除调用
    \path{clipboard_delete_item}；清空全部记录必须先通过面板内确认框，再调用
    \path{clipboard_clear_all}，避免误操作。

  \item \textbf{运行时更新设置。}
    用户保存隐私设置后，\path{clipboard_update_settings} 会先归一化并持久化新值，
    再依次停止旧监听线程和清理线程，更新共享 \path{RwLock}。若功能仍启用，则重新设置数据库容量
    并启动新的监听与清理任务；若关闭，则释放数据库状态。最后同步更新系统托盘中的启用状态，
    保证面板设置、后台任务和托盘显示一致。
\end{enumerate}
\end{sloppypar}

\subsubsection{关键代码}

\textbf{1. 监听回调与双重去重}

\begin{lstlisting}[language=Rust, caption=剪贴板变化监听]
impl ClipboardHandler for HistoryHandler {
    fn on_clipboard_change(&mut self) {
        if let Ok(text) = self.clipboard.get_text() {
            let hash = ClipboardStore::hash_content(&text);
            if consume_expected_hash(&self.expected_hash, &hash) {
                self.last_hash = Some(hash);
                return;
            }
            if self.last_hash.as_deref() == Some(&hash) { return; }
            self.last_hash = Some(hash);
            let source = current_application();
            process_text(&text, source.as_deref(),
                &self.store, &self.settings, &self.expected_hash);
        }
    }
}
\end{lstlisting}

监听器取得文本后首先消费 \texttt{expected\_hash}，避免用户从历史面板复制内容时再次入库；
随后用 \texttt{last\_hash} 过滤系统可能重复发送的变化通知。只有真正的新内容才进入隐私过滤和存储流程。

\textbf{2. SQLite 哈希去重与时间刷新}

\begin{lstlisting}[language=Rust, caption=剪贴板历史持久化]
let hash = Self::hash_content(text);
let now = Utc::now().to_rfc3339();
let conn = self.lock()?;
conn.execute(
    "INSERT INTO clipboard_history
       (content, content_hash, created_at, source_app)
     VALUES (?1, ?2, ?3, ?4)
     ON CONFLICT(content_hash) DO UPDATE SET
       content = excluded.content,
       created_at = excluded.created_at,
       source_app = COALESCE(excluded.source_app,
                             clipboard_history.source_app)",
    params![text, hash, now, source_app],
)?;
\end{lstlisting}

数据库在 \texttt{content\_hash} 上建立唯一索引。再次复制相同文本时不创建重复行，而是更新时间和
来源应用，使该记录回到列表顶部。SQLite WAL 模式还可以降低监听线程写入与面板查询之间的相互阻塞。

\textbf{3. 多关键词全文搜索}

\begin{lstlisting}[language=Rust, caption=剪贴板历史搜索条件构造]
let keywords: Vec<String> = query.split_whitespace()
    .map(str::trim)
    .filter(|part| !part.is_empty())
    .map(|part| escape_like_pattern(&part.to_lowercase()))
    .collect();

for index in 0..keywords.len() {
    if index > 0 { sql.push_str(" AND "); }
    sql.push_str(&format!(
        "LOWER(content) LIKE ?{} ESCAPE '\\'", index + 1));
}
\end{lstlisting}

搜索词按空白拆分，多个关键词必须同时命中；\texttt{\%}、\texttt{\_} 和反斜杠会先被转义，
查询值再通过参数绑定传入，从而避免用户输入改变 SQL 语义。

\subsubsection{效果截图}

\setcounter{figure}{0}
\begingroup
\setlength{\intextsep}{6pt plus 1pt minus 1pt}
\setlength{\textfloatsep}{6pt plus 1pt minus 1pt}
\captionsetup[figure]{skip=4pt}

\begin{figure}[H]
  \centering
  \includegraphics[width=0.82\textwidth]{images/clipboard-01.png}
  \caption{剪贴板管理功能总览}
\end{figure}

\begin{figure}[H]
  \centering
  \begin{minipage}[t]{0.40\textwidth}
    \centering
    \includegraphics[width=\linewidth]{images/clipboard-02.png}
    \\[0.2cm]
    \small（a）功能测试说明
  \end{minipage}
  \hfill
  \begin{minipage}[t]{0.40\textwidth}
    \centering
    \includegraphics[width=\linewidth]{images/clipboard-03.png}
    \\[0.2cm]
    \small（b）剪贴板历史列表
  \end{minipage}
  \caption{功能说明与剪贴板历史面板}
\end{figure}

\begin{figure}[H]
  \centering
  \includegraphics[width=0.58\textwidth]{images/clipboard-11.png}
  \caption{外部程序复制的历史记录}
\end{figure}

\begin{figure}[H]
  \centering
  \begin{minipage}[t]{0.40\textwidth}
    \centering
    \includegraphics[width=\linewidth]{images/clipboard-04.png}
    \\[0.2cm]
    \small（a）隐私过滤设置
  \end{minipage}
  \hfill
  \begin{minipage}[t]{0.40\textwidth}
    \centering
    \includegraphics[width=\linewidth]{images/clipboard-05.png}
    \\[0.2cm]
    \small（b）历史记录搜索
  \end{minipage}
  \caption{隐私过滤与历史记录搜索}
\end{figure}

\begin{figure}[H]
  \centering
  \includegraphics[width=0.78\textwidth]{images/clipboard-06.png}
  \caption{长文本展开效果}
\end{figure}

\begin{figure}[H]
  \centering
  \includegraphics[width=0.78\textwidth]{images/clipboard-07.png}
  \caption{一键复制成功提示}
\end{figure}

\begin{figure}[H]
  \centering
  \includegraphics[width=0.38\textwidth]{images/clipboard-08.png}
  \caption{清空全部剪贴板历史确认框}
\end{figure}

\begin{figure}[H]
  \centering
  \includegraphics[width=0.78\textwidth]{images/clipboard-09.png}
  \caption{全局搜索快捷键入口}
\end{figure}

\begin{figure}[H]
  \centering
  \includegraphics[width=0.78\textwidth]{images/clipboard-10.png}
  \caption{关键词搜索结果展示}
\end{figure}

\endgroup

\subsubsection{遇到的问题与解决方案}

\begin{enumerate}[label=(\arabic*)]
  \item \textbf{不同桌面系统的剪贴板机制不统一}
    \\→ 解决：Windows、macOS 和 X11 优先使用 \texttt{clipboard\_rs} 提供的原生事件监听；
    纯 Wayland 无法沿用 X11 事件接口，因此使用 \texttt{wl\_clipboard\_rs} 每 300 ms 读取一次，并通过原子停止标志安全结束线程。

  \item \textbf{从历史面板复制时会触发监听器，造成记录再次写入}
    \\→ 解决：写入系统剪贴板前保存目标文本的 SHA-256 到 \texttt{expected\_hash}；
    下一次监听回调命中该摘要后立即消费并跳过，1 秒内未触发事件则自动清除标记。

  \item \textbf{重复通知和重复内容导致历史列表膨胀}
    \\→ 解决：内存中的 \texttt{last\_hash} 过滤连续通知，数据库中的唯一哈希索引与 UPSERT
    负责持久化去重；重复复制只刷新时间，不新增记录。

  \item \textbf{剪贴板可能包含密码、银行卡号等敏感信息}
    \\→ 解决：存储前增加可配置过滤链。疑似密码按长度、空白、字母和数字组合判断；
    银行卡候选值必须满足 13--19 位并通过 Luhn 校验；同时支持按来源应用忽略记录。

  \item \textbf{历史记录持续增长会影响磁盘占用和查询性能}
    \\→ 解决：默认最多保存 2000 条并保留 30 天。写入后立即执行容量检查，后台每小时执行一次清理；
    超限时按最早时间以 500 条为批次删除，减少频繁的小事务。

  \item \textbf{全局快捷键可能重复触发或与应用唤醒快捷键冲突}
    \\→ 解决：注册时检测冲突并让剪贴板快捷键优先；事件处理增加 300 ms 防抖。
    面板采用居中、带标题栏的独立 Webview 窗口，标题栏 × 关闭按钮隐藏窗口而非销毁，
    重复按快捷键或点击设置面板按钮可重新打开。
\end{enumerate}

% ---------- 4.4 D ----------

\subsection{设置面板与数据管理模块（成员 D）}

\subsubsection{设计思路}

设置面板是整个 Pake Plus 的"控制中心"，负责统一管理所有模块的配置。核心设计原则是\textbf{各模块独立运行，配置集中管理}——即使面板未打开，各模块以默认配置正常运行；用户通过面板修改后即时写入本地 JSON，各模块下次读取时自动生效。

面板采用\textbf{纯 JavaScript DOM 操作}构建，无任何框架依赖。所有 UI 元素通过封装的 \texttt{el() / card() / row() / tg()} 工厂函数动态生成，支持中英文双语切换。Rust 后端负责 JSON 读写、ZIP 打包解压、系统信息采集和配置校验，通过 Tauri IPC 与前端双向通信。

前端代码（\texttt{custom.js}）以 IIFE 形式注入主窗口，在页面加载完成后通过 \texttt{setTimeout} 延迟创建左下角的齿轮入口按钮。点击按钮时动态创建包含 6 个标签页的右侧抽屉面板，各标签页通过 \texttt{buildXxx(body)} 函数按需渲染。配置数据存储在 \texttt{pake-settings.json} 中，保存前自动轮转保留最近 5 个历史版本。

\subsubsection{核心流程}

\textbf{1. 面板入口与生命周期。}
应用启动 800ms 后，\texttt{custom.js} 在页面左下角创建一个白色圆角卡片按钮（\texttt{position:fixed;bottom:24px;left:24px}），包含蓝色渐变齿轮图标和"Pake Plus"文字，hover 时边框高亮。用户点击按钮、右键托盘选择 Settings、或按 Ctrl+Shift+, 均可打开面板。面板为 \texttt{position:fixed;top:0;right:0;width:400px} 的固定定位抽屉，初始状态通过 \texttt{transform:translateX(100\%)} 隐藏在屏幕右侧，打开时 CSS transition 滑入。面板包含左侧 120px 导航栏和右侧内容区，6 个标签页通过 \texttt{buildXxx(body)} 函数按需渲染。

\textbf{2. 配置加载与表单填充。}
面板打开后调用 \texttt{get\_settings} IPC 命令加载 \texttt{AppSettings} 结构体。该结构聚合四个子模块的配置（adblock / cache / clipboard / general），每个子模块独立实现 \texttt{ModuleSettings} trait 提供 \texttt{validate() / summary() / stats()} 三个方法。加载完成后调用 \texttt{fillAll()} 遍历各标签页的 DOM 元素，通过 \texttt{tglSet()} 设置开关状态、\texttt{value} 填充选择框和文本区。

\textbf{3. 表单校验与保存。}
用户修改任何表单元素后 \texttt{dirty} 标志置为 true。点击 Save 时，\texttt{doSave()} 收集各标签页表单值构造完整 \texttt{AppSettings} 对象，调用 \texttt{validate\_settings} 校验参数合法性（如缓存范围 50-1000MB、剪贴板记录数在允许列表内等）。校验通过后调用 \texttt{save\_settings} 写入磁盘，同时触发广告拦截规则同步——后端比较新旧自定义规则，调用 \texttt{add\_custom\_rule / remove\_custom\_rule} 更新引擎。

\textbf{4. 数据导出。}
用户点击"导出"或输入路径后点"浏览"使用原生保存对话框。JS 调用 \texttt{export\_data} IPC（可选 \texttt{savePath} 参数），Rust 递归扫描 \texttt{AppData} 目录含子目录，生成 \texttt{manifest.json} 记录文件清单和元数据，用 \texttt{zip} crate 打包为 \texttt{pake-data-YYYYMMDD.zip}。

\textbf{5. 数据导入。}
用户指定 ZIP 路径或使用原生文件选择器（通过 \texttt{rfd} crate 的 \texttt{pick\_zip\_file} 命令），JS 调用 \texttt{preview\_import} 展示内容摘要，用户确认后调用 \texttt{import\_data}。Rust 解压到临时目录，按 manifest 路径映射逐文件恢复，现有文件先备份为 \texttt{.bak}。

\textbf{6. 系统诊断。}
About 标签调用 \texttt{get\_diagnostics} IPC。Rust 使用 \texttt{built} crate（编译期注入）获取版本、Git 哈希、编译时间等；使用 \texttt{sysinfo} crate 采集 OS、CPU、内存、磁盘信息。点击"复制诊断报告"通过 \texttt{arboard} crate 写入系统剪贴板。

\textbf{7. 版本备份与恢复。}
\texttt{write\_settings} 在保存前执行备份轮转：当前文件 → v1，旧 v1→v2，...，v4→v5，始终保留 5 个版本。启动时 \texttt{run\_health\_check} 检测配置是否可解析，若损坏自动从最近备份恢复。

\subsubsection{关键代码}

\textbf{1. Rust 枚举类型（类型安全）}

\begin{lstlisting}[language=Rust, caption=Theme 和 Language 枚举]
#[derive(Clone, Debug, Serialize, Deserialize, PartialEq)]
#[serde(rename_all = "lowercase")]
pub enum Theme {
    Light,
    Dark,
    System,
}

#[derive(Clone, Debug, Serialize, Deserialize, PartialEq)]
#[serde(rename_all = "lowercase")]
pub enum Language {
    #[serde(alias = "zh-CN")]
    Zh,
    En,
}
\end{lstlisting}

\textbf{2. ModuleSettings Trait（多态）}

\begin{lstlisting}[language=Rust, caption=ModuleSettings trait]
pub trait ModuleSettings {
    fn validate(&self) -> Vec<String>;      // 校验输入
    fn summary(&self) -> String;            // 一行摘要
    fn stats(&self) -> serde_json::Value;   // 统计 JSON
}

// 每个模块独立实现
impl ModuleSettings for AdblockSettings { ... }
impl ModuleSettings for CacheSettings { ... }
impl ModuleSettings for ClipboardSettings { ... }
impl ModuleSettings for GeneralSettings { ... }
\end{lstlisting}

\textbf{3. 配置校验}

\begin{lstlisting}[language=Rust, caption=设置校验逻辑]
#[command]
pub fn validate_settings(settings: AppSettings) -> Result<(), Vec<String>> {
    let mut errors: Vec<String> = Vec::new();
    errors.extend(settings.adblock.validate());    // 规则格式
    errors.extend(settings.cache.validate());      // 缓存范围 50-1000
    errors.extend(settings.clipboard.validate());  // 记录数/天数合法性
    errors.extend(settings.general.validate());    // 枚举类型编译期保证
    if errors.is_empty() { Ok(()) } else { Err(errors) }
}
\end{lstlisting}

\textbf{4. 配置校验、备份轮转与规则同步}

\begin{lstlisting}[language=Rust, caption=配置校验、备份轮转与广告规则同步]
#[command]
pub fn validate_settings(settings: AppSettings) -> Result<(), Vec<String>> {
    let mut errors = Vec::new();
    errors.extend(settings.adblock.validate());
    errors.extend(settings.cache.validate());
    errors.extend(settings.clipboard.validate());
    errors.extend(settings.general.validate());
    if errors.is_empty() { Ok(()) } else { Err(errors) }
}

const MAX_BACKUPS: u32 = 5;
fn write_settings(app: &AppHandle, s: &AppSettings) -> Result<(), String> {
    if path.exists() { /* v1→v2, ..., v4→v5 rotation */ }
    fs::write(&path, &serde_json::to_string_pretty(s).unwrap())?;
    // Sync adblock custom rules to running engine
    if let Some(st) = app.try_state::<crate::adblock::AdblockState>() {
        let cur = st.engine.custom_rules();
        if cur != s.adblock.custom_rules {
            for r in &cur { st.engine.remove_custom_rule(r); }
            for r in &s.adblock.custom_rules { st.engine.add_custom_rule(r); }
        }
    }
    Ok(())
}
\end{lstlisting}

\textbf{5. 原生文件对话框}

\begin{lstlisting}[language=Rust, caption=rfd 文件选择 IPC]
#[command] pub fn pick_save_path() -> Result<String, String> {
    rfd::FileDialog::new().set_title("选择导出位置")
        .save_file().map(|p| p.to_string_lossy().to_string())
        .ok_or_else(|| "cancelled".into())
}
#[command] pub fn pick_zip_file() -> Result<String, String> {
    rfd::FileDialog::new().set_title("选择数据文件")
        .add_filter("ZIP", &["zip"]).pick_file()
        .map(|p| p.to_string_lossy().to_string())
        .ok_or_else(|| "cancelled".into())
}
\end{lstlisting}

\subsubsection{效果截图}

% TODO: D 插入实际运行截图
\begin{itemize}
  \item 左下角 Pake Plus 入口按钮（白色卡片 + 蓝色边框 + 齿轮图标）
  \item 一般标签页——暗色主题开关、语言下拉、恢复默认按钮
  \item 广告拦截标签页——ON/OFF 开关 + 标签、✏ 编辑自定义规则折叠面板
  \item 剪贴板标签页——开关、记录数选择、保留天数、忽略短文本
  \item 数据管理标签页——导出/导入按钮、路径输入框、浏览按钮
  \item 关于标签页——诊断信息（应用信息 + 系统信息）、复制报告按钮
  \item 托盘菜单——Settings、Clipboard History、Copy Diagnostics
\end{itemize}

\subsubsection{遇到的问题与解决方案}

\begin{enumerate}[label=(\arabic*)]
  \item \textbf{Tauri initializtion\_script 中 document.body 为 null}
    \\→ 解决：将 DOM 操作包裹在 \texttt{DOMContentLoaded} 事件中，并增加即时检查兜底

  \item \textbf{切换标签页后保存导致其他标签页设置丢失}
    \\→ 解决：\texttt{doSave} 函数先检查对应元素是否存在于 DOM，不存在则沿用当前 S 状态的值

  \item \textbf{built crate 未启用 git2 feature 导致 GIT\_COMMIT\_HASH 不可用}
    \\→ 解决：在 Cargo.toml 中添加 \texttt{built = \{ version = "0.7", features = ["git2"] \}}

  \item \textbf{旧设置文件导致新默认值不生效}
    \\→ 解决：删除旧文件（AppData/Roaming 而非 Local），启动健康检查自动从备份恢复损坏文件

  \item \textbf{Theme/Language 使用 String 类型无法编译期保证合法性}
    \\→ 解决：重构为 Rust 枚举类型，\texttt{\#[serde(rename\_all = "lowercase")]} 保证序列化兼容

  \item \textbf{Tauri v2 ACL 拦截自定义 IPC 命令}
    \\→ 解决：将所有 app 级命令统一放在 \texttt{permissions/settings.toml} 的 \texttt{commands.allow} 列表中，而非分散在多个权限文件，确保所有命令被正确授权

  \item \textbf{About 标签诊断信息不显示}
    \\→ 解决：在 \texttt{buildAbout} 末尾添加 \texttt{loadDiagnostics()} 调用，截断 Git 哈希至 8 位，添加 \texttt{word-break:break-all} 防止文字溢出

  \item \textbf{右下角旧齿轮按钮与新左下角按钮并存}
    \\→ 解决：定位到原始代码中的 FAB（Floating Action Button）\texttt{\_\_ps\_fab} 和旧 \texttt{\_\_pgb} 按钮代码，逐一从 \texttt{custom.js} 中移除
\end{enumerate}

% ================================================================
% Chapter 5
% ================================================================
\section{实验结果}

\subsection{测试数据}

\begin{table}[h]
\centering
\begin{tabular}{lccc}
\toprule
测试类型 & 测试数量 & 通过 & 结果 \\
\midrule
Rust 后端测试 & 67 & 67 & 全部通过 \\
JavaScript 测试 & 297 & 297 & 全部通过 \\
cargo check & --- & --- & 0 warnings \\
cargo build & --- & --- & 成功 \\
\bottomrule
\end{tabular}
\caption{测试结果汇总}
\end{table}

\subsection{关键测试用例}

\begin{table}[h]
\centering
\begin{tabular}{p{5cm}l}
\toprule
测试用例 & 结果 \\
\midrule
默认设置合理性检查 & 通过 \\
JSON 序列化/反序列化往返 & 通过 \\
部分 JSON 解析（缺失字段回退默认值） & 通过 \\
空 JSON 解析 & 通过 \\
广告拦截规则序列化 & 通过 \\
Theme 枚举序列化为小写 & 通过 \\
Theme 枚举反序列化（合法值 + 非法值拒绝） & 通过 \\
广告规则格式校验（合法 + 非法） & 通过 \\
缓存大小范围校验（越界 + 合法） & 通过 \\
剪贴板参数校验（非法值 + 合法值） & 通过 \\
设置整体校验命令 & 通过 \\
诊断信息字段完整性 & 通过 \\
\bottomrule
\end{tabular}
\caption{关键测试用例}
\end{table}
针对成员 C 的 CLI 参数、配置合并、面板交互与权限控制进行了定向测试，
共执行 5 个测试文件、56 个测试用例，结果全部通过。

\begin{figure}[H]
  \centering
  \includegraphics[width=0.92\textwidth]{images/powershell-a-tests.png}
  \caption{PowerShell A：成员 C 定向测试结果（56 项全部通过）}
\end{figure}

CLI 边界测试分别传入 499 和 5001，程序均在打包前拒绝参数，
验证 \texttt{--clipboard-max} 仅接受 500--5000 范围内的整数。

\begin{figure}[H]
  \centering
  \includegraphics[width=0.88\textwidth]{images/powershell-a-boundary.png}
  \caption{PowerShell A：剪贴板容量参数边界校验}
\end{figure}

PowerShell B 进一步验证运行时默认值与容量压力场景。未指定容量上限的实例在空历史状态下显示
“0 条记录 / 2000”，与默认值 2000 一致；随后通过带重试的 PowerShell 函数连续写入
500 条不同的普通中文文本，用于观察监听完整性及容量上限行为。

\begin{figure}[H]
  \centering
  \includegraphics[width=0.58\textwidth]{images/powershell-b-default.png}
  \\[0.25cm]
  \includegraphics[width=0.88\textwidth]{images/powershell-b-capacity.png}
  \caption{PowerShell B：默认容量显示与 500 条剪贴板批量写入测试}
\end{figure}


\subsection{命令行示例}

\begin{lstlisting}[caption=CLI 使用示例]
# 基本打包
pake https://github.com --name GitHub

# 启用所有 Pake Plus 功能
pake https://example.com --name MyApp \
  --block-ads --adblock-rules ./my-rules.txt \
  --cache --cache-size 500 \
  --clipboard --clipboard-max 5000
\end{lstlisting}

% ================================================================
% Chapter 6
% ================================================================
\section{开源项目参考说明}

本项目基于开源项目 \textbf{Pake}（\url{https://github.com/tw93/Pake}）进行二次开发。

\subsection{引用部分}

\begin{itemize}
  \item Tauri v2 框架及窗口管理、系统托盘、全局快捷键能力
  \item CLI 打包工具链（TypeScript + Rollup）
  \item WebView 注入机制（initialization\_script）
  \item 平台适配配置（tauri.conf.json、pake.json）
\end{itemize}

\subsection{新增部分（对比原项目）}

\begin{table}[h]
\centering
\small
\begin{tabularx}{\textwidth}{lXXc}
\toprule
模块 & 新增内容 & 核心技术 & 代码量 \\
\midrule
广告拦截   & EasyList 规则引擎、网络拦截、DOM 隐藏、自定义规则 & 请求拦截器 & --- \\
离线缓存   & HTTP 代理缓存、LRU 淘汰、离线模式、缓存统计 & tokio 异步 I/O & --- \\
剪贴板     & 系统级监听（FFI）、历史搜索、一键复用、自动清理 & Win32 FFI & --- \\
设置面板   & 可视化配置、数据导入导出、诊断采集、版本回滚、健康检查 & \texttt{ModuleSettings} trait & $\sim$660 行 \\
\bottomrule
\end{tabularx}
\caption{与 Pake 原项目的差异}
\end{table}

% ================================================================
% Chapter 7
% ================================================================
\section{总结与展望}

\subsection{项目总结}

Pake Plus 成功在 Pake 的基础上实现了四个独立功能模块：

\begin{itemize}
  \item 广告拦截引擎基于全球最主流的 EasyList 规则集，网络层 + DOM 层双重过滤
  \item 离线缓存在 HTTP 层透明代理，实现断网可浏览
  \item 系统剪贴板管理通过 Rust FFI 调用操作系统 API，跨平台兼容
  \item 可视化设置面板统一管理所有配置，支持数据一键迁移和系统诊断
\end{itemize}

项目技术选型合理：Rust 负责系统级操作（文件 I/O、ZIP 打包、系统信息采集、FFI），
JavaScript 负责用户界面，Tauri IPC 作为两者之间的桥梁。
四个模块通过共享 JSON 配置文件实现了解耦协作。

\subsection{不足与改进方向}

\begin{enumerate}[label=(\arabic*)]
  \item \textbf{错误处理}：当前大量使用 \texttt{Result<(), String>}，可通过 \texttt{thiserror} 库定义结构化错误类型
  \item \textbf{异步 I/O}：IPC 命令中文件操作使用同步 I/O，可能阻塞事件循环，可迁移至 \texttt{tokio::fs}
  \item \textbf{托盘菜单动态更新}：当前无法在运行时更新菜单项文字，需重建托盘
  \item \textbf{打包时预设配置}：CLI 尚不支持将设置写入打包产物，需增加 \texttt{---settings} 选项
  \item \textbf{测试覆盖}：JS 端缺少设置面板相关的自动化测试
\end{enumerate}

\end{document}
